\subsection{Baggrund og motivation}

Verdens energiforbrug forventes at stige markant i de kommende årtier, samtidig med at behovet for stabile energikilder med lave CO$_2$ udledninger bliver stadig mere presserende.
I den sammenhæng fremstår fusionsenergi som en attraktiv mulighed, fordi den i princippet kan levere store mængder energi med begrænset brændselsforbrug
og uden de samme langsigtede affaldsproblemer som konventionel fission \cite{iter_new_energy}.
Vejen fra fysisk mulighed til teknologisk realisering er imidlertid lang, blandt andet fordi plasmaet i en fusionsreaktor er et meget komplekst system, hvor små ændringer i lokale forhold kan føre til markante forskelle i transport, stabilitet og energitab.

I denne rapport er fokus rettet mod tokamakken, hvor et varmt, ioniseret plasma holdes på plads af stærke magnetfelter.
Særligt interessant er plasmaets randregion, altså overgangen mellem det godt indesluttede plasma og scrape off layer (SOL), hvor feltlinjerne ikke længere er lukkede.
Det er netop i denne randzone, at turbulens, filamentstrukturer og tværgående transport i høj grad bestemmer både indeslutningens effektivitet og belastningen på reaktorens vægge.
Hvis disse processer ikke forstås og beskrives tilstrækkeligt præcist, bliver det vanskeligt at optimere energibalancen i reaktoren \cite{plasma_particle_sources,num_sim_blobs}.

Randplasmaets dynamik er imidlertid vanskelig at analysere direkte, fordi fysikken beskrives gennem koblede, ikke lineære partielle differentialligninger.
I dette projekt anvendes Hot Edge SOL Electrostatic (HESEL) modellen, som er en reduceret drift fluid model udviklet til at beskrive turbulens og transport
i plasmaets rand og i SOL regionen \cite{udledning_af_hesel,num_sim_blobs}.
HESEL modellen er fysisk konsistent nok til at repræsentere centrale mekanismer i randplasmaet, men samtidig tilstrækkeligt reduceret til at kunne simuleres numerisk i praksis.
Netop derfor er den et godt udgangspunkt for at undersøge, hvordan moderne datadrevne metoder kan kombineres med klassisk plasmafysisk modellering.

Den traditionelle tilgang til HESEL og lignende modeller er numerisk simulering, eksempelvis ved finite difference diskretisering i rum og tidsintegration af det resulterende ligningssystem.
Denne tilgang er veletableret, men også beregningstung, og hvis man ønsker at sammenligne mange scenarier eller udforske store parameterrum, bliver fulde simuleringer hurtigt en flaskehals.
Problemet er derfor ikke alene at kunne simulere systemet korrekt, men også at kunne gøre det effektivt og på en måde, der bevarer den væsentlige fysik.

Her bliver Scientific Machine Learning (SciML) interessant. SciML betegner metoder, hvor maskinlæring kombineres med kendt struktur fra naturvidenskabelige modeller, eksempelvis differentialligninger og bevarelseslove.
For plasmafysik er det særligt lovende, fordi mange problemer netop er kendetegnet ved dyre simuleringer og stærke fysiske begrænsninger, som modellerne bør respektere.
En vellykket SciML model kan derfor potentielt fungere som en hurtigere surrogatmodel i udvalgte arbejdsgange.

Projektets motivation er først og fremmest metodisk: hvis SciML kan anvendes meningsfuldt på HESEL modellen,
peger det mod en mere fleksibel måde at arbejde med komplekse plasmamodeller på.
Projektet placerer sig derfor i krydsfeltet mellem anvendt matematik, numerisk analyse, plasmafysik og maskinlæring, og netop dette tværfaglige spændingsfelt gør undersøgelsen fagligt relevant.

\subsection{Formål og forskningsspørgsmål}

Formålet med rapporten er at undersøge, hvordan SciML kan anvendes til modellering af randdynamik i plasma med udgangspunkt i HESEL modellen,
og i hvilken grad sådanne modeller kan supplere eller effektivisere klassiske numeriske simuleringer.
Udgangspunktet er ikke, at maskinlæring skal erstatte fysikbaserede metoder, men at afklare, hvornår en fysikinformeret læringsmodel er et nyttigt redskab.

% På baggrund af projektplanen undersøges følgende spørgsmål:
Derfor stilles følgende forskningsspørgsmål:
\begin{enumerate}
    \item Hvordan kan SciML, særligt fysikinformerede neurale netværk, anvendes til modeller- ing af HESEL ligningssystemet, og hvordan sammenlignes løsningerne med klassiske numeriske simuleringer
    \item Hvilken betydning har det for modellens kvalitet, om man vægter de Partielle Differential ligninger (PDE)'er og data'en forskelligt, og hvordan påvirker dataopløsningen resultaterne? 
    \item I hvor høj grad kan modellen generalisere til ukendte punkter, tidsskridt og simulations skonfigurationer 
    \item Hvor meget kan modellen effektiviseres, før tabet i præcision bliver væsentligt, og hvordan bør beregningstid, hukommelsesforbrug og nøjagtighed afvejes i praksis?
\end{enumerate}

Disse spørgsmål er valgt, fordi de tilsammen dækker både den fysiske og den beregningsmæssige side af problemet. Det er ikke tilstrækkeligt, at en model blot kan reproducere enkelte løsninger; den skal også være robust, fortolkbar og anvendelig i situationer, hvor man ønsker mere end ren interpolation. Derfor handler projektet både om modelleringsevne og om videnskabelig troværdighed.

Rapporten har samtidig et læringsperspektiv. Arbejdet forudsætter forståelse af numerisk løsning af PDE systemer, grundlæggende plasmafysik i randregionen og udvikling samt evaluering af SciML modeller på dynamiske systemer. 
% Læringsmålene fungerer derfor som en bagvedliggende ramme for projektet.

Resten af rapporten er struktureret som følger. Først præsenteres den fysiske baggrund for HESEL modellen og de numeriske metoder, der anvendes til simulering af ligningssystemet. Herefter beskrives implementeringen af eksperimenterne og den måde, data genereres og behandles på. Dernæst analyseres resultaterne fra både de klassiske simuleringer og de anvendte SciML metoder med henblik på at vurdere nøjagtighed, generalisering og effektivitet. Afslutningsvis diskuteres projektets væsentligste resultater, begrænsninger og perspektiver for fremtidigt arbejde.
